\documentclass[b5paper]{jsarticle}

\setlength{\columnseprule}{0.4pt}
\setlength{\oddsidemargin}{-5mm}
\setlength{\topmargin}{-20mm}
\setlength{\textwidth}{140mm}
\setlength{\textheight}{220mm}
%\pagestyle{empty}
%\usepackage{amsmath}
\special{papersize=\the\paperwidth,\the\paperheight}



\usepackage[dvipdfmx,hiresbb]{graphicx}
\usepackage{amsmath,amssymb}
%\usepackage{fancybox}
 \usepackage{ascmac}
 \usepackage{multicol}
 \usepackage{multicol,multienum}
 \usepackage{latexsym}
  \usepackage{wrapfig}
  
\begin{document}
\newcounter{nombre} %必須
\renewcommand{\thenombre}{\arabic{nombre}} %任意
\setcounter{nombre}{0} %任意

\newcommand{\toi}[1][]{\refstepcounter{nombre}#1\fbox{\large\makebox[1ex]{\textbf{\thenombre}}} \hspace{1mm}}

\newcommand{\MARU}[1]{{\ooalign{\hfil#1\/\hfil\crcr\raise.167ex\hbox{\mathhexbox20D}}}}

\def\heiko{\mathrel{\raise.3ex\hbox{\scalebox{.7}{%
    \rotatebox[origin=c]{-7}{/}%
    \kern-.35em\rotatebox[origin=c]{-7}{/}}}}}%




\vspace*{\stretch{1}}
	\begin{flushleft}
	\HUGE\textbf{4年\ 最難関講習\ 数学\\第8章「複素数と方程式」\\確認テスト}
	\end{flushleft}

	\begin{flushleft}
	\LARGE\textbf{試験時間　30分}
	\end{flushleft}
\vspace{20mm}
 注意事項
 
 \begin{description}
\item[1]  この冊子の解答欄に解答を作成しなさい。
\item[2] 計算式だけでなく途中過程や考え方を丁寧に記述しなさい。
\end{description}



\vspace{\stretch{2}}

 \underline{\hspace{1cm}組\hspace{1cm}番\hspace{5mm}氏名\hspace{7cm}} \hspace{5mm}\underline{\hspace{2cm}点}
\newpage
\setcounter{nombre}{0} 
\vspace{5mm}
\toi%1

次の方程式を解け．
\begin{description}
\item[(1)] $x^3-21x^2+128x-180=0$
\item[(2)] $(x^2-2x)(x^2-2x-4)+4=0$
\end{description}

【解答欄】

\newpage


　


\newpage
\toi%2

実数を係数とする3次方程式$x^3+ax^2+bx+3=0$が$1+\sqrt2 i$を解にもつ．
\begin{description}
\item[(1)] $a,b$の値を求めよ．
\item[(2)]  他の2つの解を求めよ．
\item[(3)] 3つの解を$\alpha ,\beta ,\gamma$とするとき，$\alpha^5+\beta^5+\gamma^5$の値を求めよ．
\end{description}

【解答欄】

\newpage

　


\newpage
\toi%3


2次方程式$x^2+kx+k^2-2=0$の2解を$\alpha ,\beta$とする．
\begin{description}
\item[(1)] $\alpha^2+\beta^2=1$であるとき，$k$の値を求めよ．
\item[(2)] $\alpha ,\beta$がともに正になるような$k$の範囲を求めよ．
\item[(3)] $\alpha ,\beta$の少なくとも一つが正であるような$k$の範囲を求めよ．
\end{description}

【解答欄】

\newpage


　



\newpage

\end{document}